\chapter{Аналитическая часть}

\section{Синтетические даннные}

Синтетические данные – это искусственно сгенерированные компьютером данные,
обладающие теми же характеристиками и статистическими свойствами, что и
реальные данные.
Они применяются в том случае, если получение реальных данных затруднено, требует больших затрат, ставит под угрозу конфиденциальность пользователей или занимает слишком много времени.
Синтетические данные могут быть текстовыми или табличными, а
также представлять собой видео, изображения или звуки, как и реальные данные.
%https://elibrary.ru/item.asp?id=49251096

% Дальше вода

%синтетические данные очень гибкие. Использование
%таких данных позволяет создавать, распространять и исключать
%синтетические данные по своему усмотрению с целью исключения
%предвзятости ИИ. Их использование способно повысить качество данных.
%Под качеством данных будем понимать степень соответствия данных
%требованиям потребителя
%https://elibrary.ru/item.asp?id=49595169

%синтетические данные являются альтернативой
%историческим или отсутствующим данным. Общеизвестна проблема
%отсутствующих значений в данных. При работе с реальными данными
%необходимо заполнять отсутствующие значения средним
%арифметическим на определенной выборке, модой, медианой и т.д. [3, 4].
%Генерируемые алгоритмами ИИ синтетические данные в настоящее время
%в основном используются для обучения моделей машинного обучения.
%Однако они начинают появляться в разработке и тестировании
%программного обеспечения. Такое решение значительно сокращает время
%разработки и вывода продукта на рынок
%https://elibrary.ru/item.asp?id=49595169

%синтетические данные безопасны и полностью анонимны.
%Вместо того, чтобы изменять и в процессе уничтожать существующий
%набор данных, синтетические данные создаются с нуля
%https://elibrary.ru/item.asp?id=49595169


Синтетические данные разделяют на три категории:
\begin{enumerate}
	\item Полностью синтетические данные, которые не имеют с реальными
	данными почти ничего общего и генерируются на основе некоторых признаков
	реальных данных. Они сохраняют статистические свойства реальных данных и
	фактически заменяют их.
	\item Частично синтетические данные, которые заменяют только
	некоторые значения в реальных данных, если они являются конфиденциальными
	или содержат пропуски.
	\item Гибридные синтетические данные, которые генерируются с
	использованием и реальных, и полностью синтетических данных: для реальных
	данных выбираются близкие им синтетические, и они объединяются, формируя
	гибридные данные.
\end{enumerate}
%https://elibrary.ru/item.asp?id=49251096

Создание синтетических данных включает в себя различные методы и подходы, которые позволяют генерировать данные, максимально приближенные к реальным по своим характеристикам и структуре.

Одним из основных методов является статистическое моделирование, при котором данные генерируются на основе заранее определённых статистических моделей и распределений. Например, можно использовать нормальное распределение для создания данных, похожих на реальные по среднему значению и стандартному отклонению.
%https://naukaip.ru/wp-content/uploads/2024/06/K-649.pdf#page=12 
%Явная Д. Р. СИНТЕТИЧЕСКИЕ ДАННЫЕ: СОЗДАНИЕ И ИСПОЛЬЗОВАНИЕ В РАЗРАБОТКЕ //ЛУЧШАЯ ИССЛЕДОВАТЕЛЬСКАЯ РАБОТА 2024: сборник статей. – 2024. – С. 12

\subsection{Генеративыне модели}
Генеративные модели представляют собой преобразующий подход в
системах принятия решений, фиксируя сложные распределения данных и
генерируя новые экземпляры данных, которые сохраняют исходные
характеристики данных, которые они моделируют. Эта внутренняя способность
определяет их роль в структурах принятия решений, где прогнозирование
новых результатов на основе прошлых данных имеет решающее значение.
https://cyberleninka.ru/article/n/razvitie-generativnyh-modeley-v-logicheskih-sistemah-prinyatiya-resheniy/viewer



\subsection{Статические и вероятностные методы}

\section{Дифференциальная приватность}

Дифференциальная приватность обеспечивает защиту личной информации, гарантируя, что результаты анализа не позволят точно определить информацию о конкретных индивидах. 
В рамках ДП данные обрабатываются с добавлением шума, что минимизирует влияние индивидуальных записей на результаты статистических запросов.
Таким образом, как показано на рисунке~\ref{img:nirs}, даже если в набор данных добавляется или исключается информация о конкретном индивиде, это не влияет на общие выводы о группе~\cite{Wood2018DifferentialPA}.

Представим БД, содержащую медицинские записи 100 человек, из которых 20 страдают диабетом.
Злоумышленник хочет узнать, есть ли у человека $X$ диабет, и уже узнал, что 19 из остальных 99 человек в выборке являются диабетиками.
Запрашивая базу данных и получая информацию о том, что в выборке 20 диабетиков, злоумышленник делает вывод, что человек
$X$ также страдает диабетом.
Этот процесс извлечения дедуктивной информации называется дифференцированной атакой~\cite{harvard}.

Теперь представим, что доступ к данным, хранящихся в БД, предоставляется через систему запросов с дифференциальной приватностью.
Когда злоумышленник делает запрос, система возвращает истинное значение с некоторым случайным шумом.
%Выходные данные могут составлять 21, 25 или 19 в зависимости от случайного шума.
Даже если злоумышленник знает все о других 99 людях, то он не может быть уверен, есть ли у человека $X$ диабет.
Конфиденциальность человека $X$ защищена, поскольку его риск нарушения конфиденциальности остается примерно одинаковым независимо от того, входит ли он в набор данных или нет.
Это справедливо для любого отдельного человека и любого набора данных, независимо от того, насколько необычными или отличительными могут быть данные конкретного индивидуума.~\cite{harvard}
\newpage

\includeimage{nirs}{f}{h}{1\textwidth}{Принцип ДП}


\section{Типы приватности}

Параметр приватности $\epsilon$ --- это неотрицательное значение, определяющее уровень приватности, который влияет на количество добавляемого шума~\cite{harvard}.

\subsection{$\epsilon$ -- дифференцированная приватность}
Модель $\epsilon$ -- дифференцированной приватности гарантирует, что результаты анализа данных не изменятся, если одна запись будет добавлена или удалена.~\cite{main}

Механизм $M$ обеспечивает $\epsilon$ -- дифференциальную приватность, если для любых двух соседних наборов данных $D$ и $D'$ (различающихся на одну запись) выполняется условие по формуле~\ref{equ:eps-diff}:

\begin{equation}\label{equ:eps-diff}
	Pr[M(D)] \eqslantless \exp(\epsilon) \times Pr[M(D')]
\end{equation}
, где $\epsilon$ --- параметр приватности

\subsection{$(\epsilon$, $\delta)$ -- дифференцированная приватность}
Модель $(\epsilon$, $\delta)$ -- дифференцированной приватности вводит дополнительный параметр $\delta$, который ослабляет строгие гарантии приватности.
Параметр $\delta$ определяет вероятность того, что гарантия приватности будет нарушена в некоторых случаях.
Это дает гибкость в балансе между приватностью и точностью результатов.~\cite{main}

Механизм $M$ обеспечивает ($\epsilon$, $\delta$) -- дифференциальную приватность, если для любых двух соседних наборов данных $D$ и $D'$ (различающихся на одну запись) выполняется условие по формуле~\ref{equ:eps-delta-diff}

\begin{equation}\label{equ:eps-delta-diff}
	Pr[M(D)] \eqslantless \exp(\epsilon) \times Pr[M(D')] + \delta
\end{equation}
, где $\epsilon$ --- параметр приватности, $\delta$ --- параметр вероятности нарушения приватности



\section{Типы механизмов}

Параметр чувствительности~---~это максимальное изменение результата механизма при изменении одного элемента в базе данных.

Чувствительностью механизма $M$ является максимум разности значений механизма для любых двух соседних наборов данных $D$ и $D'$ (различающихся на одну запись) и вычисляется по формуле~\ref{eq:sens}:
\begin{equation}\label{eq:sens}
	\Delta M  = \max(|M(D) - M(D')|)
\end{equation}
, где $M$ --- параметр чувствительности.

\subsection{Механизм Лапласа}
Лапласов механизм добавляет к результату запроса шум, сгенерированный из распределения Лапласа~\cite{main}.
Результат Лапласова механизма вычисляется по формуле~\ref{eq:laplas}:
\begin{equation}\label{eq:laplas}
	M'(D) = M(D) + \text{Lap}(0, \frac{\Delta M}{\epsilon})
\end{equation}
где $\Delta M$ --- параметр чувствительности, $\epsilon$ --- параметр приватности

\subsection{Механизм Гаусса} 
Гауссов механизм похож на Лапласов, но шум добавляется на основе нормального (Гауссова) распределения.
К результатам запроса добавляется шум по формуле~\ref{eq:gauss-norm}:

\begin{equation}\label{eq:gauss-norm}
	M'(D) = M(D) + \mathcal{N}(0, \sigma^2)
\end{equation}
, где $\mathcal{N}(0, \sigma^2)$ --- нормальное распределение.

Параметр $\sigma$ вычисляется по формуле~\ref{eq:gauss-sigma}.
\begin{equation}\label{eq:gauss-sigma}
	\sigma = \frac{\Delta M \sqrt{2 \ln(1.25/\delta)}}{\epsilon}
\end{equation}

\subsection{Экспоненциальный механизм}
Экспоненциальный механизм работает путем добавления случайности в
процесс выбора элементов, где вероятность выбора каждого элемента зави-
сит от его "качества"(оценки полезности) и параметра приватности.

Для двух соседних наборов данных $D$ и $D'$  экспоненциальный механизм $M$ обеспечивает экспоненциальную приватность, если вероятность выбора результата $r$ зависит от его полезности $u(r, D)$, и вычисляется по формуле~\ref{eq:exp-mech}
\begin{equation}\label{eq:exp-mech}
	P[M(D)] \propto \exp(\epsilon  u(r, D))
\end{equation}
, где $\epsilon$ --- параметр приватности, $u(r, D)$ --- оценка полезности 